\documentclass[11pt,a4paper]{article}
\usepackage[italian]{babel}
\usepackage[margin=2.5cm]{geometry}
\usepackage{parskip}
\usepackage{amsmath, amssymb, amsthm}
\usepackage{xcolor}
\usepackage{listings}
\usepackage{tikz}
\usepackage[most]{tcolorbox}
\usepackage{hyperref}

% --- Riquadri con bordo nero, senza numero ---
\tcbset{nota/.style={colback=white, colframe=black, boxrule=0.5pt,
  sharp corners}}
\NewTColorBox{definizione}{}{nota,
  before upper={\textbf{Definizione.}\ }}
\NewTColorBox{osservazione}{}{nota,
  before upper={\textbf{Osservazione.}\ }}
\theoremstyle{definition}
\newtheorem*{esempio}{Esempio}
\newtheorem*{esercizio}{Esercizio}

% --- Codice C ---
% Uso: \begin{codice} ... \end{codice}; \begin{risultato} ... \end{risultato}
%      \lstinline|printf("ciao");| per il codice dentro al testo
\lstdefinestyle{c}{
  language=C,
  basicstyle=\ttfamily\small,
  keywordstyle=\color{blue}\bfseries,
  commentstyle=\color{gray}\itshape,
  stringstyle=\color{red!70!black},
  numbers=left,
  numberstyle=\tiny\color{gray},
  numbersep=8pt,
  columns=flexible,
  keepspaces=true,
  breaklines=true,
  showstringspaces=false,
  tabsize=4,
  morekeywords={include, define},
  literate={à}{{\`a}}1 {è}{{\`e}}1 {é}{{\'e}}1 {ì}{{\`i}}1
           {ò}{{\`o}}1 {ù}{{\`u}}1
}
\lstset{style=c}
\newtcblisting{codice}{listing only, nota, left=6mm,
  listing options={style=c}}
\newtcblisting{risultato}{listing only, colback=black!5, colframe=black,
  boxrule=0.5pt, sharp corners,
  listing options={basicstyle=\ttfamily\small, numbers=none, language={}}}

% --- Diagrammi di flusso (simboli standard) ---
\usetikzlibrary{shapes.geometric, arrows.meta, positioning}
\tikzset{
  terminale/.style = {ellipse, draw, minimum width=2.5cm, minimum height=0.9cm},
  io/.style        = {trapezium, trapezium left angle=70, trapezium right angle=110,
                      draw, minimum width=2.5cm, minimum height=0.9cm},
  processo/.style  = {rectangle, draw, minimum width=2.5cm, minimum height=0.9cm},
  decisione/.style = {diamond, draw, aspect=2, minimum width=2.5cm},
  freccia/.style   = {thick, -{Stealth}}
}

\title{Scope e binding}
\author{Marco Cerbella}
\date{Lezione 7 -- 6 ottobre 2026}

\begin{document}
\maketitle

\section{Scope degli identificatori}

\begin{definizione}
Lo \emph{scope} (ambito) di un identificatore è la parte del programma in cui
l'identificatore è significativo, cioè in cui può essere ``visto''. Il tipo di
scope dipende da dove l'identificatore viene dichiarato.
\end{definizione}

\begin{itemize}
    \item \textbf{File scope}: l'identificatore è dichiarato fuori da tutti i
          blocchi e dalle liste di parametri (cioè fuori dalle funzioni). Si
          può usare dalla dichiarazione fino alla fine del file.
    \item \textbf{Block scope}: l'identificatore è dichiarato dentro un blocco.
          Si può usare dalla dichiarazione fino alla fine del blocco più piccolo
          che la contiene (spesso, ma non sempre, il corpo di una funzione).
\end{itemize}

\begin{esempio}
\begin{codice}
int p;                   // file scope

int fun2(void) {         // fun2: file scope
    int r = 3;           // r, s: block scope
    int s = r;
    while (r != s) {
        r--;
    }
    return r;
}
\end{codice}
\end{esempio}

\begin{osservazione}
Non si possono avere due variabili con lo stesso nome nello stesso scope:
\begin{codice}
int main() {
    int a = 5;
    float a = 6.5;   // errore: a già dichiarata
}
\end{codice}
\end{osservazione}

\section{Shadowing}

Si può riusare un identificatore in una dichiarazione \emph{annidata} nello
scope esistente; la nuova dichiarazione deve avere block scope, in un blocco
contenuto in quello esterno. In questo caso la dichiarazione interna
\emph{nasconde} (\emph{shadowing}) quella esterna, che non è visibile nello
scope interno.

\begin{esempio}
\begin{codice}
int main() {
    int x = 1, y = 2, z = 3;
    printf(" x = %d, y = %d, z = %d \n", x, y, z);
    {
        int x = 10;
        float y = 20;
        printf(" x = %d, y = %f, z = %d \n", x, y, z);
        {
            int z = 100;
            printf(" x = %d, y = %f, z = %d \n", x, y, z);
        }
    }
    return 0;
}
\end{codice}
\begin{risultato}
x = 1, y = 2, z = 3
x = 10, y = 20.000000, z = 3
x = 10, y = 20.000000, z = 100
\end{risultato}
\end{esempio}

\begin{esempio}
Qui \lstinline|x++| modifica la \lstinline|x| esterna (non è ridichiarata),
mentre \lstinline|y++| modifica la \lstinline|y| interna; uscendo dal blocco
interno la \lstinline|y| torna a valere 20.
\begin{codice}
int main() {
    {
        int x = 10, y = 20;
        {
            printf("x = %d, y = %d\n", x, y);
            {
                int y = 40;
                x++;
                y++;
                printf("x = %d, y = %d\n", x, y);
            }
            printf("x = %d, y = %d\n", x, y);
        }
    }
    return 0;
}
\end{codice}
\begin{risultato}
x = 10, y = 20
x = 11, y = 41
x = 11, y = 20
\end{risultato}
\end{esempio}

\begin{esempio}
Fuori dal blocco in cui è dichiarata, una variabile non è accessibile:
\begin{codice}
int main() {
    {
        int x = 10;
    }
    {
        printf("%d", x);    // errore: x non visibile qui
    }
    return 0;
}
\end{codice}
\begin{risultato}
error: 'x' undeclared (first use in this function)
\end{risultato}
\end{esempio}

\section{Binding statico}

\begin{definizione}
Un \emph{binding} è l'associazione tra un nome e l'oggetto che esso denota. Il
\emph{binding time} è il momento in cui il binding viene creato (più in
generale, il momento in cui viene presa una decisione di implementazione). La
regione testuale del programma in cui un binding è attivo è il suo scope.
\end{definizione}

In C lo scope di un binding è determinato \emph{staticamente}, cioè a tempo di
compilazione (\emph{static scoping}), come si vede dal testo del programma e non
dall'ordine delle chiamate a run-time. Altri linguaggi usano lo scoping
\emph{dinamico}: ad esempio in Perl \lstinline|my| definisce una variabile
locale a scope statico, \lstinline|local| una a scope dinamico.

\section{Esercizio}

\begin{esercizio}
Qual è l'output del seguente programma?
\begin{codice}
#include <stdio.h>

int f1(void);
int f2(int x, int a);

int a;

int main() {
    int a, b, c;

    a = 7;
    b = f1();
    c = f2(a, b);
    printf("%d %d %d\n", a, b, c);
}

int f1(void) {
    a = 12;
    printf("%d ", a);
    return (a + 5);
}

int f2(int x, int a) {
    printf("%d ", a);
    return (x * a);
}
\end{codice}
\end{esercizio}

\emph{Soluzione.}
\begin{itemize}
    \item In \lstinline|main| la \lstinline|a| locale (che vale 7) nasconde la
          \lstinline|a| globale.
    \item \lstinline|f1| non dichiara una propria \lstinline|a|: per binding
          statico modifica quella \emph{globale}, la pone a 12, stampa
          \texttt{12} e restituisce 17, quindi $b = 17$.
    \item In \lstinline|f2(7, 17)| il parametro \lstinline|a| vale 17: stampa
          \texttt{17} e restituisce $7 \cdot 17 = 119$, quindi $c = 119$.
    \item Infine \lstinline|main| stampa \texttt{7 17 119}.
\end{itemize}
\begin{risultato}
12 17 7 17 119
\end{risultato}
\end{document}
